\chapter{Como Ler as Fontes sem Transformar Referências em Decoração}
\label{ap:fontes}

A bibliografia identifica fontes concretas desta edição. Trabalhos históricos continuam relevantes: o artigo original do Transformer explica uma arquitetura, não todos os modelos atuais. Um paper que mediu ganhos de CoT em certas tarefas não prova ganhos em qualquer sistema de atendimento.

\section*{Um protocolo de leitura}
Localiza a afirmação a verificar. Confirma título, autores, versão e data. Lê o método e identifica modelos, dados, comparações e métricas. Procura limitações e verifica se a conclusão fica dentro do que foi medido. Guarda página ou secção e distingue a frase dos autores da tua interpretação.

Se só tens o resumo, confirma apenas o tema e os resultados que ele declara; não afirmes ter auditado detalhes metodológicos. Se uma página oficial exige acesso que não tens, regista a limitação e procura outra fonte primária. Não preenches a lacuna com uma citação inventada.

\section*{Percursos de aprofundamento}
Para redes neuronais, \emph{Deep Learning} oferece uma base de referência~\citep{goodfellow2016}. Para linguagem, \emph{Speech and Language Processing} organiza fundamentos e aplicações; a edição em linha evolui e deve ser identificada pela data consultada~\citep{jurafsky}.

Lê Brown, Wei, Kojima e Wang com atenção às tarefas e aos modelos. Confronta-os com estudos de explicações infiéis e autocorreção. Para sistemas, começa por RAG e ReAct e compara as arquiteturas com o teu problema. Protocolos e APIs exigem documentação da versão concreta.

OWASP oferece uma taxonomia de falhas; as especificações permitem implementar interfaces. No direito, o texto em vigor tem uma função diferente de um tutorial. O registo \texttt{docs/FONTES\_VERIFICADAS.md} documenta o alcance da consulta e o que continua a exigir revisão humana.
